\subsection{半双线性映射。}

本节中，除非另有明确说明，我们用 A 与 B 表示两个环，用 E 表示左 A-模，用 F 表示左 B-模；我们用 \(b \to b^{\mathrm{J}} (b \in \mathrm{B})\) 表示 B 的一个反自同构，即 B 到其自身上的一个双射，满足 \((b + c)^{\mathrm{J}} = b^{\mathrm{J}} + c^{\mathrm{J}}\) 与 \((bc)^{\mathrm{J}} = c^{\mathrm{J}}b^{\mathrm{J}}\) （对 B 中的一切 \(b, c\)）；我们用 \(\mathbf{J}'\) 表示 \(\mathbf{J}^{-1}\) 。我们用 G 表示一个 \((\mathbf{A}, \mathbf{B})\) -双模（第 1 小节）。

定义 4. --- 称映射 \(\Phi\) （从 E \(\times\) F 到 G 中）为关于 J 的右半双线性映射，如果它满足条件 (1)、(2)、(3)（定义 1，第 1 小节）以及

\begin{enumerate}
\def\labelenumi{(\arabic{enumi})}
\setcounter{enumi}{6}
\tightlist
\item
  \(\Phi(x, by) = \Phi(x, y). b^J\) 对一切 \(x \in \mathbf{E}, y \in \mathbf{F}\) 与 \(b \in \mathbf{B}\) 。
\end{enumerate}

若 J 是恒同映射（这要求 B 是交换的），我们就重新得到双线性映射的概念。

设 \((e_i)_{i\in I}\) 与 \((f_k)_{k\in K}\) 是 E 与 F 的元素的两个族，并设 \((a_i)_{i\in I}\) 与 \((b_k)_{k\in K}\) 是 A 与 B 的元素，其中除有限个外均为零。我们有

\[
\Phi (\sum_ {i} a _ {i} e _ {i}, \sum_ {k} b _ {k} f _ {k}) = \sum_ {i, k} a _ {i} \Phi (e _ {i}, f _ {k}) b _ {k} ^ {\mathrm{J}}.\tag{8}
\]

与双线性映射的情形一样，元素 \(\Phi(e_i, f_k)\) 以唯一的方式确定 \(\Phi\) （当 \((e_i)\) 与 \((f_k)\) 是生成系时），而当 \((e_i)\) 与 \((f_k)\) 是 E 与 F 的基时可任意选取；当 \((e_i)\) 与 \((f_k)\) 是有限基时，称 \((\Phi(e_i, f_k))\) 为 \(\Phi\) 关于这些基的矩阵。

与双线性映射一样，在 E × F 到 G 中的（关于 J 的）右半双线性映射的集合上，可定义 A 与 B 的中心上的一个双模结构。半双线性映射的原像概念，用与双线性映射相同的公式来定义。此外我们将看到，半双线性映射的研究可以化为双线性映射的研究。

定义 5. --- 设 B 是一个环，F 是左（相应地右）B-模，J 是 B 的一个反自同构。用 F \(^{j}\) 表示具有与 F 相同的底加法群、且外合成律为 \((b, y) \to b^{J'}y\) （相应地 \((b, y) \to yb^{J'}\) ）\((b \in B, y \in F, J' = J^{-1})\) 的右（相应地左）B-模。

按定义 5 的记号，从 \(\mathbf{F}^{\mathrm{j}}\) 到一个右（相应地左）B-模 H 中的线性映射，因而就等同于 F 到 H 的一个 Z-线性映射 \(u\) ，它满足

\[
u (b y) = u (y) b ^ {J} \quad (\text { resp. } u (y b) = b ^ {J} u (y)) \quad (b \in B, y \in F).
\]

映射 \(u\) 是 F 到 H 的、关于 J 的半线性映射（第 II 章，附录 I，第 1 小节），如果把 J 看作 B 的反环 B \(^{0}\) 到 B 上的一个同构，并把 F 看作一个右（相应地左）B \(^{0}\) -模。

类似地，当 F 是左 B-模时，E × F 到 G 中的（关于 J 的）右半双线性映射 \(\Phi\) 等同于 E × F \(^{J}\) 到 G 中的一个双线性映射；若后者右退化（相应地左退化、非退化），则称 \(\Phi\) 右退化（相应地左退化、非退化）。

注记. --- 设 A 与 B 为两个环，\(J_{1}\) 为 A 的一个反自同构，M 为右 A-模，N 为右 B-模，G 为 (A, B)-双模。称映射 \(\Phi\) （从 M \(\times\) N 到 G 中）为关于 J \(_{1}\) 的左半双线性映射，如果它是 Z-双线性的且满足

\[
\Phi (x a, y b) = a ^ {\mathrm{J} _ {1}} \Phi (x, y) b \quad (x \in \mathrm{M}, y \in \mathrm{N}, a \in \mathrm{A}, b \in \mathrm{B}).\tag{9}
\]

这样的映射等同于 \(M^{J_{1}} \times N\) 到 G 中的一个双线性映射。我们常把为右半双线性映射给出的定义与性质向左半双线性映射的转述留给读者；当我们说到半双线性映射（不加指明）时，它总是右半双线性映射。

